\subsection{MAXIMAL TORI AND COVERING OF HOMOMORPHISMS}

Let G be a connected compact Lie group, T a maximal torus of G. Consider the derived group \(\mathrm{D}(\mathrm{G})\) of G and its universal covering \(\tilde{\mathrm{D}}(\mathrm{G})\) ; let \(p: \tilde{\mathrm{D}}(\mathrm{G}) \to \mathrm{G}\) be the composite of the canonical morphisms \(\tilde{\mathrm{D}}(\mathrm{G}) \to \mathrm{D}(\mathrm{G})\) and \(\mathrm{D}(\mathrm{G}) \to \mathrm{G}\) . Then \(\tilde{\mathrm{D}}(\mathrm{G})\) is a connected compact Lie group ( \(\S1\) , no. 4, Cor. 2 of Prop. 4); moreover, the inverse image \(\tilde{T}\) of T under p is a maximal torus of \(\tilde{\mathrm{D}}(\mathrm{G})\) (no. 3, Prop. 1).

Lemma 2. Let H be a Lie group, \(f_{\mathrm{T}}: \mathrm{T} \to \mathrm{H}\) and \(\tilde{f}: \tilde{\mathrm{D}}(\mathrm{G}) \to \mathrm{H}\) morphisms of Lie groups such that \(f_{\mathrm{T}}(p(t)) = \tilde{f}(t)\) for all \(t \in \tilde{\mathrm{T}}\) . There exists a unique morphism of Lie groups \(f: \mathrm{G} \to \mathrm{H}\) such that \(f \circ p = \tilde{f}\) and such that the restriction of \(f\) to \(\mathrm{T}\) is \(f_{\mathrm{T}}\) .

Put \(Z = \mathrm{C}(\mathrm{G})_{0}\) ; by §1, no. 4, Cor. 1 of Prop. 4, the morphism of Lie groups \(g: Z \times \tilde{\mathrm{D}}(\mathrm{G}) \to \mathrm{G}\) such that \(g(z, x) = z^{-1} p(x)\) is a covering; its kernel consists of the pairs \((z, x)\) such that \(p(x) = z\) , for which \(x \in p^{-1}(\mathrm{Z}) \subset \tilde{\mathrm{T}}\) . Since the morphism \((z, x) \mapsto f_{\mathrm{T}}(z^{-1}) \tilde{f}(x)\) from \(Z \times \tilde{\mathrm{D}}(\mathrm{G})\) to H maps \(\operatorname{Ker} g\) to \(\{e\}\) , there exists a morphism f from G to H such that \(f \circ p = \tilde{f}\) and \(f(z) = f_{\mathrm{T}}(z)\) for \(z \in Z\) . But we also have \(f(t) = f_{\mathrm{T}}(t)\) for \(t \in p(\tilde{\mathrm{T}})\) ; since \(T = Z.p(\tilde{\mathrm{T}})\) , the restriction of f to T is indeed \(f_{T}\) .

PROPOSITION 8. Let G be a connected compact Lie group, T a maximal torus of G, H a Lie group and \(\varphi : \mathrm{L}(\mathrm{G}) \to \mathrm{L}(\mathrm{H})\) a homomorphism of Lie algebras. There exists a morphism of Lie groups \(f: \mathrm{G} \to \mathrm{H}\) such that \(\mathrm{L}(f) = \varphi\) if and only if there exists a morphism of Lie groups \(f_{\mathrm{T}}: \mathrm{T} \to \mathrm{H}\) such that \(\mathrm{L}(f_{\mathrm{T}}) = \varphi | \mathrm{L}(\mathrm{T})\) ; then \(f_{\mathrm{T}} = f | \mathrm{T}\) .

If \(f: \mathrm{G} \to \mathrm{H}\) is a morphism of Lie groups such that \(\mathrm{L}(f) = \varphi\) , the restriction \(f_{\mathrm{T}}\) of \(f\) to \(\mathrm{T}\) is the unique morphism from \(\mathrm{T}\) to \(\mathrm{H}\) such that \(\mathrm{L}(f_{\mathrm{T}}) = \varphi | \mathrm{L}(\mathrm{T})\) . Conversely, let \(f_{\mathrm{T}}: \mathrm{T} \to \mathrm{H}\) be a morphism of Lie groups such that \(\mathrm{L}(f_{\mathrm{T}}) = \varphi | \mathrm{L}(\mathrm{T})\) . Let \(\tilde{\mathrm{D}}(\mathrm{G})\) and \(p\) be as above; the map \(\mathrm{L}(p)\) induces an isomorphism from \(\mathrm{L}(\tilde{\mathrm{D}}(\mathrm{G}))\) to the derived algebra \(\mathfrak{b}\) of \(\mathrm{L}(\mathrm{G})\) . There exists a morphism of Lie groups \(\tilde{f}: \tilde{\mathrm{D}}(\mathrm{G}) \to \mathrm{H}\) such that \(\mathrm{L}(\tilde{f}) = (\varphi | \mathfrak{b}) \circ \mathrm{L}(p)\) (Chap. III, §6, no. 1, Th. 1). The morphisms \(t \mapsto \tilde{f}(t)\) and \(t \mapsto f_{\mathrm{T}}(p(t))\) from \(\tilde{\mathrm{T}}\) to \(\mathrm{H}\) induce the same homomorphism of Lie algebras, and hence coincide. Applying Lemma 2, we deduce the existence of a morphism \(f: \mathrm{G} \to \mathrm{H}\) such that \(\mathrm{L}(f)\) and \(\varphi\) coincide on \(\mathrm{L}(\mathrm{T})\) and \(\mathfrak{b}\) . Since \(\mathrm{L}(\mathrm{G}) = \mathfrak{b} + \mathrm{L}(\mathrm{T})\) , we have \(\mathrm{L}(f) = \varphi\) .

PROPOSITION 9. Let G be a connected compact Lie group, T a maximal torus of G, H a Lie group and \(f: G \to H\) a morphism. Then f is injective if and only if its restriction to T is injective.

Indeed, by Th. 2 (no. 2) the normal subgroup \(\mathrm{Ker}f\) of G reduces to the identity element if and only if its intersection with T reduces to the identity element.

